\section{Experiments}





\begin{table}[t]
\centering
\caption{\textbf{Comparison to state-of-the-art methods on Waymo (val).} 
PSNR and SSIM are reported. Full: requires dense scene annotations.
Camera: requires camera intrinsics and extrinsics.}
{\footnotesize\setlength{\tabcolsep}{.8mm}
\resizebox{\columnwidth}{!}{
\begin{tabular}{lccccc}
\toprule[1.2pt]                      
\multirow{2}{*}{\textbf{Methods}} & \multirow{2}{*}{\textbf{Supervision}}  & \multicolumn{2}{c}{\textbf{Dynamic-only}}  & \multicolumn{2}{c}{\textbf{Full image}}\\ 
\cmidrule(lr){3-4} \cmidrule(lr){5-6}
& & $\text{PSNR}\uparrow$ & $\text{SSIM}\uparrow$ & $\text{PSNR}\uparrow$ & $\text{SSIM}\uparrow$\\ 
\midrule
\textbf{Per-scene optimization} \\
3DGS \cite{kerbl20233d} & Full &17.13 &0.267  &25.13 &0.741 \\
DeformableGS \cite{yang2024deformable} & Full &17.10 &0.266  &\textbf{25.29} &\textbf{0.761} \\
\midrule
\textbf{Feed-forward model} \\
GS-LRM\cite{zhang2024gs} & Camera &20.02 &0.520 &25.18 &0.753 \\
STORM\cite{yang2024storm} & Camera &\textbf{21.26} &\textbf{0.535} &25.03 &0.750 \\
DynamicVGGT & Image-only &18.07 &0.376 &24.07 &0.676 \\
\bottomrule[1.2pt]
\end{tabular}
}
}
\label{tab:psnr}
\end{table}


% \begin{table}[t]
% \centering
% \caption[Valori medi]{Comparison of flow estimation on Waymo(val)}
% {\footnotesize\setlength{\tabcolsep}{.8mm}
% {\begin{tabular}{lcccc}
% \toprule[1.2pt]                      
% {           \textbf{Methods}}                & 
%           $ \text{EPE3D(m)} \downarrow $     & 
%           $ \text{$Acc_{5}$} \uparrow  $     & 
%           $ \text{$Acc_{10}$} \uparrow $     & 
%           $ \theta\text{(rad)} \downarrow $     
% \\ 
% \midrule
% NSFP\cite{li2021neural}            &0.698 &42.17 &54.26 &0.919    \\
% STORM\cite{yang2024storm}           &\textbf{0.276} &\textbf{81.12} &\textbf{85.61} &\textbf{0.658}   \\
% DynamicVGGT     &- &-  & -  & -             \\
% \bottomrule[1.2pt]
% \end{tabular}
% }
% }
% \label{tab:flow}
% \end{table}

% \begin{table}[t]
% \centering
% \caption[Valori medi]{Monocular and MVS depth estimation.}
% {\footnotesize\setlength{\tabcolsep}{.8mm}
% \resizebox{\columnwidth}{!}{
% {\begin{tabular}{lcccccc}
% \toprule[1.2pt]                      
% \multirow{2}{*}{\textbf{Methods}} & \multicolumn{2}{c}{\textbf{{KITTI(Mono)}}}  & \multicolumn{2}{c}{\textbf{{NYU-v2(Mono)}}} & \multicolumn{2}{c}{\textbf{{KITTI(MVS)}}}
% \\ 
% \cmidrule(lr){2-3} \cmidrule(lr){4-5} \cmidrule(lr){6-7}
%           & $ \text{Abs Rel} \downarrow $   & 
%           $ \delta<1.25 \uparrow $          & 
%           $\text{Abs Rel} \downarrow  $     & 
%           $ \delta<1.25 \uparrow $       & 
%           $\text{Abs Rel} \downarrow  $     & 
%           $ \delta<1.25 \uparrow $            
% \\ 
% \midrule
% DUSt3R\cite{wang2024dust3r}          & 0.109             & 0.873              & 0.081            & 0.909      &0.143 &0.814     \\
% MASt3R\cite{leroy2024grounding}          & 0.077             & \textbf{0.948}  & 0.110            & 0.865      &0.115 &0.848     \\
% MonST3R         & 0.098             & 0.895              & 0.094            & 0.887      &0.107 &0.884      \\
% VGGT            & 0.082             & 0.938              & 0.059            & 0.951      &0.062 &0.969     \\
% StreamVGGT      & 0.082             & 0.947              & 0.057            & 0.959     &0.173 &0.721       \\
% % $\pi^3$         & \textbf{0.060}    & 0.971              & 0.054            & 0.956           \\
% DynamicVGGT     & \textbf{0.070} & 0.940  & 0.064            & 0.943     &\textbf{0.051} &\textbf{0.976}      \\

% \bottomrule[1.2pt]
% \end{tabular}
% }
% }
% }
% \label{tab:mono-depth-eval}
% \end{table}

% \begin{table*}[t]
% \centering
% \caption[Valori medi]{\textbf{Ablation study.}We evaluate the ablated variants of DynamicVGGT by recording their point map estimation on KITTI and Waymo(val) datasets. KITTI uses monocular input with every 3 consecutive frames per camera.
% Waymo uses 3 frames (stride 4) from FRONT, SIDE LEFT, and SIDE RIGHT cameras, totaling 9 images per group.}
% {\footnotesize\setlength{\tabcolsep}{.8mm}
% {\begin{tabular}{lcccccccccccccccccc}
% \toprule[1.2pt]                      
% \multirow{3}{*}{\textbf{Methods}} & \multicolumn{6}{c}{\textbf{{KITTI(Mono)}}}   & \multicolumn{6}{c}{\textbf{{Waymo(3cam)}}}
% \\ 
% \cmidrule(lr){2-7} \cmidrule(lr){8-13} 
% & \multicolumn{3}{c}{\textbf{{Mean}}} & \multicolumn{3}{c}{\textbf{{Med.}}} &
% \multicolumn{3}{c}{\textbf{{Mean}}} & \multicolumn{3}{c}{\textbf{{Med.}}}
% \\ 
% \cmidrule(lr){2-4} \cmidrule(lr){5-7} \cmidrule(lr){8-10} \cmidrule(lr){11-13} 
%           & $ \text{Acc.} \downarrow $     & 
%           $ \text{Comp.} \downarrow  $     &
%           $ \text{NC} \uparrow  $     & 
%           $ \text{Acc.} \downarrow   $     & 
%           $ \text{Comp.} \downarrow    $     & 
%           $ \text{NC} \uparrow   $     & 
%           $ \text{Acc.} \downarrow   $     & 
%           $ \text{Comp.} \downarrow    $     &
%           $ \text{NC} \uparrow  $     & 
%           $ \text{Acc.} \downarrow $     & 
%           $ \text{Comp.} \downarrow  $     &
%           $ \text{NC} \uparrow  $     &   
% \\ 
% \midrule
% Baseline & 1.489 & 0.690 & 0.918 & 1.329 & 0.535 & \textbf{0.971}    & 4.635 & 2.667 & 0.561 & 2.634 & 1.734 & 0.590 \\
%  \quad + TA \& FPH(stage1) &0.927 & 0.600 & 0.915 & 0.857  & 0.474 &0.932 & 4.330 & 2.939 & 0.561 & 2.224  & 1.649 & 0.593 \\
%  \quad + DGSH(stage2)  &\textbf{0.901} &\textbf{0.584} &\textbf{0.939} &\textbf{0.733} & \textbf{0.464} &\underline{0.963}   &\textbf{4.021} &2.390 &\textbf{0.562} &\textbf{1.971} &\textbf{1.564} &\underline{0.603} \\
% \bottomrule[1.2pt]
% \end{tabular}
% }
% }
% \label{tab:ab-study}
% \end{table*}


This section compares our method to state-of-the-art approaches across multiple tasks to show its effectiveness.




\subsection{Implementation Details}

% By default, we use $L = 12$ layers of MTA, resulting in a total of approximately 1.4B parameters.
% DynamicVGGT is initialized from the pretrained VGGT weights, and about 800M parameters (excluding frozen modules) are fine-tuned in two stages.
% In Stage 1, we train for 10 epochs using the AdamW optimizer with a hybrid schedule—linear warm-up for the first 0.5 epoch followed by cosine decay—peaking at a learning rate of $1\times10^{-6}$.
% In Stage 2, we fine-tune for an additional 50 epochs with the Gaussian head enabled, using the same schedule but a higher peak learning rate of $5\times10^{-5}$.
% All input frames, depth maps, and point maps are resized such that the longer image side does not exceed 518 pixels. The variable temporal step $\Delta t$, randomly sampled within a preset range from 1 to 3. We employ a dynamic batch sizing approach similar to VGGT, processing 18 images
% per GPU. For the training objective, we set $\lambda_1$ = 0.01, $\lambda_2$=0.1, and $\lambda_3$=0.01.
% Training is conducted on 32 $\times$ H200 GPUs, taking about 6 hours for Stage 1 and a day for Stage 2.
% Further details of the training settings and dataset configurations are provided in the Appendix.

% 删减版本
We use $L=12$ MTA layers, resulting in approximately 1.4B parameters. 
DynamicVGGT is initialized from pretrained VGGT weights, with about 800M parameters (excluding frozen modules) fine-tuned in two stages.

In Stage~1, we train for 10 epochs using AdamW with a hybrid learning rate schedule: a linear warm-up for the first 0.5 epoch followed by cosine decay, with a peak learning rate of $1\times10^{-6}$. 
In Stage~2, we further fine-tune the model for 50 epochs with the Gaussian head enabled, using the same schedule but a higher peak learning rate of $5\times10^{-5}$.

All input frames, depth maps, and point maps are resized so that the longer image side does not exceed 518 pixels. 
The temporal offset $\delta$ is randomly sampled from 1 to 3. 
We adopt a dynamic batch sizing strategy similar to VGGT, processing 18 images per batch.
For the training objective, we set 
$\lambda_{\mathrm{temp}} = 0.01$, 
$\lambda_{\mathrm{gs}} = \lambda_{\mathrm{dis}} = 0.1$, 
and $\lambda_{\mathrm{flow}} = 0.01$. Further details of the training settings and dataset configurations are provided in the Appendix.

\subsection{Training Data}
The model is trained on a collection of dynamic autonomous driving datasets, including Waymo Open Dataset~\cite{waymo}, Virtual KITTI~\cite{virtualkitti}, and MVS-Synth~\cite{mvssynth}.
In the first stage, DynamicVGGT is trained on Virtual KITTI and MVS-Synth to learn robust geometric priors and temporal consistency under controlled synthetic settings.
In the second stage, the model is fine-tuned with the Dynamic 3D Gaussian Splatting module on Waymo and Virtual KITTI to enhance dynamic geometry refinement and appearance consistency in real driving environments.
Evaluation is conducted on the Waymo validation set and the KITTI~\cite{kitti17} dataset to assess reconstruction quality.

% More implementation and evaluation details are provided in the Appendix.
\subsection{Point Map Reconstruction}
We evaluate point map reconstruction on the KITTI and Waymo datasets, as shown in Table~\ref{tab:point-estimation}.
We report Accuracy (Acc.), Completion (Comp.), and Normal Consistency (NC). On the KITTI dataset, which uses monocular input with three consecutive frames per sequence, DynamicVGGT achieves the best results across most metrics, obtaining an accuracy of 0.901 and a normal consistency of 0.939.
It consistently outperforms both VGGT~\cite{wang2025vggt} and StreamVGGT~\cite{zhuo2025streaming}, demonstrating its effectiveness in capturing dynamic geometry and maintaining temporal consistency in monocular sequences.

On the Waymo dataset, which provides synchronized multi-view images from three cameras with a frame stride of four, our model generalizes well to large-scale dynamic driving scenes.
It achieves an accuracy of 4.021 and a normal consistency of 0.603.
These results confirm that the proposed dynamic formulation effectively enhances cross-view consistency and scene completeness, even under challenging real-world motion and illumination variations, highlighting the scalability of our feed-forward framework for dynamic 4D perception.



% 放附录
% \begin{table*}[t]
% \centering
% \caption[Valori medi]{\textbf{Point Map Reconstruction on KITTI, Waymo and ETH3D.} KITTI uses monocular input with every 3 consecutive frames per camera. Waymo uses 3 frames (stride 4) from FRONT, SIDE\_LEFT, and SIDE\_RIGHT cameras, totaling 9 images per group. }
% {\footnotesize\setlength{\tabcolsep}{.8mm}
% {\begin{tabular}{lcccccccccccccccccc}
% \toprule[1.2pt]                      
% \multirow{3}{*}{\textbf{Methods}} & \multicolumn{6}{c}{\textbf{{KITTI (monocular)}}}   & \multicolumn{6}{c}{\textbf{{Waymo (3 cameras)}}} & \multicolumn{6}{c}{\textbf{{ETH3D}}}
% \\ 
% \cmidrule(lr){2-7} \cmidrule(lr){8-13} \cmidrule(lr){14-19}
% & \multicolumn{3}{c}{\textbf{{Mean}}} & \multicolumn{3}{c}{\textbf{{Med}}} &
% \multicolumn{3}{c}{\textbf{{Mean}}} & \multicolumn{3}{c}{\textbf{{Med}}} &
% \multicolumn{3}{c}{\textbf{{Mean}}} & \multicolumn{3}{c}{\textbf{{Med}}} 
% \\ 
% \cmidrule(lr){2-4} \cmidrule(lr){5-7} \cmidrule(lr){8-10} \cmidrule(lr){11-13} \cmidrule(lr){14-16} \cmidrule(lr){17-19}
%           & $ \text{Acc.} \downarrow $     & 
%           $ \text{Comp.} \downarrow  $     &
%           $ \text{NC} \uparrow  $     & 
%           $ \text{Acc.} \downarrow   $     & 
%           $ \text{Comp.} \uparrow    $     & 
%           $ \text{NC} \uparrow   $     & 
%           $ \text{Acc.} \downarrow   $     & 
%           $ \text{Comp.} \uparrow    $     &
%           $ \text{NC} \uparrow  $     & 
%           $ \text{Acc.} \downarrow $     & 
%           $ \text{Comp.} \downarrow  $     &
%           $ \text{NC} \uparrow  $     & 
%           $ \text{Acc.} \downarrow   $     & 
%           $ \text{Comp.} \uparrow    $     & 
%           $ \text{NC} \uparrow   $     & 
%           $ \text{Acc.} \downarrow   $     & 
%           $ \text{Comp.} \uparrow    $     &
%           $ \text{NC} \uparrow  $     
% \\ 
% \midrule

% zsy annotated: 测试结果 at 2025.11.12。可以任选一个进去。
% % 7scenes
% VGGT:       & 0.079 & 0.085 & 0.792 & 0.034 & 0.037 & 0.891
% StreamVGGT: & 0.125 & 0.104 & 0.749 & 0.050 & 0.038 & 0.862
% % DTU
% VGGT:       & 0.698 & 0.530 & 0.952 & 0.468 & 0.363 & 0.991
% StreamVGGT: & 2.414 & 1.213 & 0.865 & 1.676 & 0.654 & 0.971
% % ETH3D
% VGGT        & 0.737 & 0.770 & 0.815 & 0.491 & 0.375 & 0.935
% StreamVGGT: & 0.339 & 0.473 & 0.868 & 0.244 & 0.233 & 0.959
% % KITTI
% VGGT:       & 1.489 & 0.690 & 0.918 & 1.329 & 0.535 & 0.971
% StreamVGGT: & 1.078 & 0.495 & 0.899 & 0.867 & 0.390 & 0.949
% % NRGBD
% VGGT:       & 0.079 & 0.082 & 0.909 & 0.020 & 0.023 & 0.990
% StreamVGGT: & 0.084 & 0.079 & 0.861 & 0.044 & 0.039 & 0.986
% Waymo
% VGGT:       & 4.635 & 2.667 & 0.561 & 2.634 & 1.734 & 0.590
% StreamVGGT: & 4.598 & 2.626 & 0.564 & 2.567 & 1.789 & 0.592
% VGGT            & 1.489 & 0.690 & 0.918 & 1.329 & 0.535 & 0.971    & 4.635 & 2.667 & 0.561 & 2.634 & 1.734 & 0.590    & 0.737 & 0.770 & 0.815 & 0.491 & 0.375 & 0.935  \\
% StreamVGGT      & 1.078 & 0.495 & 0.899 & 0.867 & 0.390 & 0.949    & 4.598 & 2.626 & 0.564 & 2.567 & 1.789 & 0.592    & 0.339 & 0.473 & 0.868 & 0.244 & 0.233 & 0.959  \\
% Anysplat        &- &- &- &-  &- &-                   \\
% % DynamicVGGT     &\textbf{0.4774} &1.2497  &0.8213 &\textbf{0.2700} &\textbf{0.3343}  &\textbf{0.9396} \\
% DynamicVGGT     &\textbf{0.901} &\underline{0.584} &\textbf{0.939} &\textbf{0.733} &\underline{0.464} &\underline{0.963}                  \\

% \bottomrule[1.2pt]
% \end{tabular}
% }
% }
% \label{tab:point-estimation}
% \end{table*}


% 不比相机，首先没创新，其次相机的内参变化太小，相机放附录里
% \subsubsection{Camera Pose Estimation}
% We assess point map reconstruction quality across both scene-level
% and object-level datasets: 7-Scenes, NRGBD, and
% DTU. We use multi-view images with fixed sequence-id mappings from for fair comparison, reporting Accuracy (Acc.) and Completion (Comp.) metrics in
% Tab. 1. These results clearly demonstrate our
% model’s ability to effectively leverage prior information for better reconstruction.

\subsection{4D scene reconstruction}

We further evaluate DynamicVGGT on 4D scene reconstruction using the Waymo validation set, as summarized in Table~\ref{tab:psnr}.
We compare two categories of methods: per-scene optimization and feed-forward models.
DynamicVGGT achieves a PSNR of 18.07 and an SSIM of 0.376 on dynamic regions, and reaches 24.07 and 0.676 on the full-frame evaluation.
Although methods such as STORM obtain higher scores with multi-camera inputs and geometric priors, achieving 21.26 in PSNR and 0.535 in SSIM, DynamicVGGT delivers competitive results using only monocular images without relying on camera parameters or scene-specific optimization.
These results demonstrate that DynamicVGGT effectively reconstructs dynamic 4D scenes with strong temporal consistency and high visual fidelity through purely image-based self-supervision.

% \subsubsection{Flow estimation}
% \begin{table}[t]
% \centering
% \caption[Valori medi]{Comparison of flow estimation on Waymo(val)}
% {\footnotesize\setlength{\tabcolsep}{.8mm}
% {\begin{tabular}{lcccc}
% \toprule[1.2pt]                      
% {           \textbf{Methods}}                & 
%           $ \text{EPE3D(m)} \downarrow $     & 
%           $ \text{$Acc_{5}$} \uparrow  $     & 
%           $ \text{$Acc_{10}$} \uparrow $     & 
%           $ \theta\text{(rad)} \downarrow $     
% \\ 
% \midrule
% NSFP\cite{li2021neural}            &0.698 &42.17 &54.26 &0.919    \\
% STORM\cite{yang2024storm}           &\textbf{0.276} &\textbf{81.12} &\textbf{85.61} &\textbf{0.658}   \\
% DynamicVGGT     &- &-  & -  & -             \\
% \bottomrule[1.2pt]
% \end{tabular}
% }
% }
% \label{tab:flow}
% \end{table}
%  A key capability of DynamicVGGT is scene motion estimation, which we evaluate on the Waymo Open Dataset. Following \cite{li2021neural}, we report standard metrics: 3D End-Point Error (EPE3D), Acc5, Acc10, angular error ($\theta$ err). We compare against NSFP\cite{li2021neural} and STORM\cite{yang2024storm} in Table \ref{tab:flow}. 



% \begin{table}[h]
% \centering
% \caption[Valori medi]{Comparison of flow estimation on Waymo(val)}
% {\footnotesize\setlength{\tabcolsep}{.8mm}
% {\begin{tabular}{lcccc}
% \toprule[1.2pt]                      
% {           \textbf{Methods}}                & 
%           $ \text{EPE3D(m)} \downarrow $     & 
%           $ \text{$Acc_{5}$} \uparrow  $     & 
%           $ \text{$Acc_{10}$} \uparrow $     & 
%           $ \theta\text{(rad)} \downarrow $     
% \\ 
% \midrule
% NSFP\cite{li2021neural}            &0.698 &42.17 &54.26 &0.919    \\
% STORM\cite{yang2024storm}           &\textbf{0.276} &\textbf{81.12} &\textbf{85.61} &\textbf{0.658}   \\
% DynamicVGGT     &- &-  & -  & -             \\
% \bottomrule[1.2pt]
% \end{tabular}
% }
% }
% \label{tab:flow}
% \end{table}

\subsection{Monocular and MVS depth estimation}
\begin{table}[t]
\centering
\caption[Valori medi]{Monocular and MVS depth estimation.}
{\footnotesize\setlength{\tabcolsep}{.8mm}
\resizebox{\columnwidth}{!}{
{\begin{tabular}{lcccccc}
\toprule[1.2pt]                      
\multirow{2}{*}{\textbf{Methods}} & \multicolumn{2}{c}{\textbf{{KITTI(Mono)}}}  & \multicolumn{2}{c}{\textbf{{NYU-v2(Mono)}}} & \multicolumn{2}{c}{\textbf{{KITTI(MVS)}}}
\\ 
\cmidrule(lr){2-3} \cmidrule(lr){4-5} \cmidrule(lr){6-7}
          & $ \text{Abs Rel} \downarrow $   & 
          $ \delta<1.25 \uparrow $          & 
          $\text{Abs Rel} \downarrow  $     & 
          $ \delta<1.25 \uparrow $       & 
          $\text{Abs Rel} \downarrow  $     & 
          $ \delta<1.25 \uparrow $            
\\ 
\midrule
DUSt3R\cite{wang2024dust3r}          & 0.109             & 0.873              & 0.081            & 0.909      &0.143 &0.814     \\
MASt3R\cite{leroy2024grounding}          & 0.077             & \textbf{0.948}  & 0.110            & 0.865      &0.115 &0.848     \\
MonST3R \cite{zhang2024monst3r}         & 0.098             & 0.895              & 0.094            & 0.887      &0.107 &0.884      \\
VGGT\cite{wang2025vggt}            & 0.082             & 0.938              & 0.059            & 0.951      &0.062 &0.969     \\
StreamVGGT \cite{zhuo2025streaming}     & 0.082             & 0.947              & \textbf{0.057}            & \textbf{0.959}     &0.173 &0.721       \\
% $\pi^3$         & \textbf{0.060}    & 0.971              & 0.054            & 0.956           \\
DynamicVGGT     & \textbf{0.070} & 0.940  & 0.064            & 0.943     &\textbf{0.051} &\textbf{0.976}      \\

\bottomrule[1.2pt]
\end{tabular}
}
}
}
\label{tab:mono-depth-eval}
\end{table}
\begin{table*}[t]
\centering
\caption[Valori medi]{\textbf{Ablation study.} We evaluate ablated variants of DynamicVGGT on point map estimation over KITTI and Waymo (val). KITTI uses monocular input with three consecutive frames, while Waymo uses 3 frames (stride 4) from the FRONT, SIDE LEFT, and SIDE RIGHT cameras, yielding 9 images per sample. Metrics include Accuracy (Acc.), Completeness (Comp.), Normal Consistency (NC).}
{\footnotesize\setlength{\tabcolsep}{.8mm}
{\begin{tabular}{lcccccccccccccccccc}
\toprule[1.2pt]                      
\multirow{3}{*}{\textbf{Methods}} & \multicolumn{6}{c}{\textbf{{KITTI(Mono)}}}   & \multicolumn{6}{c}{\textbf{{Waymo(3cam)}}}
\\ 
\cmidrule(lr){2-7} \cmidrule(lr){8-13} 
& \multicolumn{3}{c}{\textbf{{Mean}}} & \multicolumn{3}{c}{\textbf{{Med.}}} &
\multicolumn{3}{c}{\textbf{{Mean}}} & \multicolumn{3}{c}{\textbf{{Med.}}}
\\ 
\cmidrule(lr){2-4} \cmidrule(lr){5-7} \cmidrule(lr){8-10} \cmidrule(lr){11-13} 
          & $ \text{Acc.} \downarrow $     & 
          $ \text{Comp.} \downarrow  $     &
          $ \text{NC} \uparrow  $     & 
          $ \text{Acc.} \downarrow   $     & 
          $ \text{Comp.} \downarrow    $     & 
          $ \text{NC} \uparrow   $     & 
          $ \text{Acc.} \downarrow   $     & 
          $ \text{Comp.} \downarrow    $     &
          $ \text{NC} \uparrow  $     & 
          $ \text{Acc.} \downarrow $     & 
          $ \text{Comp.} \downarrow  $     &
          $ \text{NC} \uparrow  $     &   
\\ 
\midrule
Baseline & 1.489 & 0.690 & 0.918 & 1.329 & 0.535 & \textbf{0.971}    & 4.635 & 2.667 & 0.561 & 2.634 & 1.734 & 0.590 \\
 \quad + TA \& FPH(stage1) &0.927 & 0.600 & 0.915 & 0.857  & 0.474 &0.932 & 4.330 & 2.939 & 0.561 & 2.224  & 1.649 & 0.593 \\
 \quad + DGSHead(stage2)  &\textbf{0.901} &\textbf{0.584} &\textbf{0.939} &\textbf{0.733} & \textbf{0.464} &\underline{0.963}   &\textbf{4.021} &2.390 &\textbf{0.562} &\textbf{1.971} &\textbf{1.564} &\underline{0.603} \\
\bottomrule[1.2pt]
\end{tabular}
}
}
\label{tab:ab-study}
\end{table*}

\begin{figure}[t]
\centering
  \includegraphics[width=1\columnwidth]{figs/pointmap_restructionv4.pdf}
  \caption{\textbf{Point map reconstruction.} 
DynamicVGGT reconstructs denser, smoother, and more geometrically consistent point maps than VGGT, maintaining temporal coherence even under large viewpoint or scene changes. Zoom in for better view.}
  \label{fig:point_est}
\end{figure}

\begin{figure*}[t]
\centering
  \includegraphics[width=\textwidth]{figs/3dgsv2.pdf}
  \caption{\textbf{Scene Reconstruction and Novel View Synthesis.} Given input frames 0, 2, and 4, DynamicVGGT reconstructs the corresponding scenes and synthesizes novel views for the next frame. For both KITTI and Waymo, multi-view inputs are used; for Waymo, we visualize the front-camera results due to limited view overlap. The model achieves high-quality reconstruction and realistic novel view generation across dynamic driving scenes.}
  \label{fig:3dgs}
\end{figure*}

We evaluate the monocular and multi-view stereo depth estimation performance of DynamicVGGT on three benchmarks: KITTI and NYU-v2~\cite{SilbermanECCV12}, as shown in Table~\ref{tab:mono-depth-eval}. 
We report two standard metrics: Absolute Relative Error (Abs Rel) and accuracy under the $\delta<1.25$ threshold. 

On monocular KITTI, DynamicVGGT achieves an Abs Rel of 0.070, outperforming all baselines. 
On NYU-v2, it obtains an Abs Rel of 0.064 and 94.3\% accuracy under $\delta<1.25$, demonstrating strong generalization from outdoor to indoor scenes. 
Under the multi-view stereo setting, DynamicVGGT achieves the best overall results with an Abs Rel of 0.051 and 97.6\% accuracy, surpassing VGGT and StreamVGGT by a clear margin.

% 放可视化，量化放附录
\subsection{Ablution study}

We conduct an ablation study on the KITTI monocular dataset and the Waymo three-camera dataset to evaluate the contribution of each proposed component, as summarized in Table~\ref{tab:ab-study}.
Starting from the vanilla VGGT baseline, adding temporal attention and the future point prediction head improves accuracy from 1.489 to 0.927 and completeness from 0.690 to 0.600 on KITTI, demonstrating the benefit of temporal modeling in capturing dynamic geometry.

Introducing the Dynamic 3D Gaussian Splatting Head further enhances accuracy to 0.901 and normal consistency to 0.939, producing smoother and more complete reconstructions.
On the Waymo dataset, our full model achieves the best overall results with an accuracy error of 4.021 and a normal consistency of 0.603, confirming that the combination of motion-aware temporal attention and dynamic geometry refinement substantially improves temporal coherence and reconstruction quality.



\subsection{Visualization}

\noindent \textbf{Point Map Reconstruction.}  
We provide qualitative comparisons of point map reconstruction results in Fig.~\ref{fig:point_est}, covering three configurations: single-frame reconstruction, short-term multi-frame reconstruction, and long-range temporal reconstruction.
Across all settings, DynamicVGGT consistently outperforms VGGT in both geometric completeness and temporal consistency.
Even in the single-frame case, our model produces denser and smoother geometry with more accurate structural details, whereas VGGT tends to generate distorted geometry when the viewpoint changes.

When extended to multi-frame and long-sequence inputs, DynamicVGGT demonstrates superior robustness in capturing point-level motion trajectories and maintaining consistent global geometry, particularly in challenging scenarios such as downhill roads and open intersections in autonomous driving.
This highlights the model’s ability to recover fine-grained 3D structures and maintain stable temporal coherence under large scene variations.

\noindent \textbf{Scene Reconstruction and Novel View Synthesis.}  
Figure~\ref{fig:3dgs} presents qualitative results of scene reconstruction and novel view synthesis on the real-world KITTI and Waymo datasets.
Given input frames ${0, 2, 4}$, DynamicVGGT reconstructs the corresponding scenes and synthesizes novel views for the subsequent frame, showcasing its ability to model temporal dynamics from purely image-based inputs.
For both datasets, multi-view images are used as input; for Waymo, we visualize the front-camera results due to the limited overlap between different camera views.

Our method faithfully reconstructs dynamic scenes with moving vehicles and illumination changes, while maintaining consistent global geometry and realistic appearance across time.
Notably, DynamicVGGT achieves temporally coherent novel view synthesis, highlighting its effectiveness as a unified feed-forward framework for dynamic 4D perception in autonomous driving.






% \begin{table*}[t]
% \centering
% \caption[Valori medi]{\textbf{Point Map Reconstruction on KITTI and Waymo(val).} KITTI uses monocular input with every 3 consecutive frames per camera. Waymo uses 3 frames (stride 4) from FRONT, SIDE\_LEFT, and SIDE\_RIGHT cameras, totaling 9 images per group. }
% {\footnotesize\setlength{\tabcolsep}{.8mm}
% {\begin{tabular}{lcccccccccccccccccc}
% \toprule[1.2pt]                      
% \multirow{3}{*}{\textbf{Methods}} & \multicolumn{6}{c}{\textbf{{KITTI (monocular)}}}   & \multicolumn{6}{c}{\textbf{{Waymo (3 cameras)}}}
% \\ 
% \cmidrule(lr){2-7} \cmidrule(lr){8-13} \cmidrule(lr){14-19}
% & \multicolumn{3}{c}{\textbf{{Mean}}} & \multicolumn{3}{c}{\textbf{{Med.}}} &
% \multicolumn{3}{c}{\textbf{{Mean}}} & \multicolumn{3}{c}{\textbf{{Med.}}} 
% \\ 
% \cmidrule(lr){2-4} \cmidrule(lr){5-7} \cmidrule(lr){8-10} \cmidrule(lr){11-13}
%           & $ \text{Acc.} \downarrow $     & 
%           $ \text{Comp.} \downarrow  $     &
%           $ \text{NC} \uparrow  $     & 
%           $ \text{Acc.} \downarrow   $     & 
%           $ \text{Comp.} \downarrow    $     & 
%           $ \text{NC} \uparrow   $     & 
%           $ \text{Acc.} \downarrow   $     & 
%           $ \text{Comp.} \downarrow    $     &
%           $ \text{NC} \uparrow  $     & 
%           $ \text{Acc.} \downarrow $     & 
%           $ \text{Comp.} \downarrow  $     &
%           $ \text{NC} \uparrow  $     &   
% \\ 
% \midrule

% VGGT\cite{wang2025vggt}            & 1.489 & 0.690 & 0.918 & 1.329 & \textbf{0.535} & \textbf{0.971}    & 4.635 & 2.667 & 0.561 & 2.634 & 1.734 & 0.590  \\
% StreamVGGT\cite{zhuo2025streaming}      & 1.078 & \textbf{0.495} & 0.899 & 0.867 & 0.390 & 0.949    & 4.598 & 2.626 & 0.564 & 2.567 & 1.789 & 0.592  \\
% % DynamicVGGT     &\textbf{0.4774} &1.2497  &0.8213 &\textbf{0.2700} &\textbf{0.3343}  &\textbf{0.9396} \\
% DynamicVGGT     &\textbf{0.901} &\underline{0.584} &\textbf{0.939} &\textbf{0.733} &\underline{0.464} &\underline{0.963}   &\textbf{4.021} &\textbf{2.390} &\textbf{0.562} &\textbf{1.971} &\underline{1.564} &\underline{0.603}                \\

% \bottomrule[1.2pt]
% \end{tabular}
% }
% }
% \label{tab:point-estimation}
% \end{table*}